\documentclass{article}
\usepackage[utf8]{inputenc}
\usepackage{amsmath}
\usepackage{amssymb}
\usepackage{graphicx}
\usepackage{tabularx}
\usepackage{booktabs}
\usepackage{caption}
\usepackage[english]{babel}
\usepackage{setspace}
\usepackage{geometry}
\usepackage{float}
\usepackage[justification=centering]{caption}

% APA style formatting
\geometry{a4paper, margin=1in}
\setstretch{1.5}
\captionsetup[table]{labelfont=bf, justification=raggedright, singlelinecheck=false}
\captionsetup[figure]{font=small, labelfont=bf, justification=raggedright, singlelinecheck=false}

\title{ Machine Learning for Customer Churn Prediction}
\author{Nelson  \\ Javier  \\ Nicolas Castro Rivera }
\date{}

\begin{document}

\maketitle

\begin{abstract}

Customer Churn is a problem about customer discontinue a company's services or products. Machine learning techniques can be applied with purpose of predict if a customer 
has probability of be churn, this information can help the company to take decisions for not lost clients. This work search explore different machine learning techniques
to compare and get the best approach to solve this problem. 

\end{abstract}

\section{Introduction}
Customer Churn refers to clients that potentialy can discontinue a service or a product of a company. The main context for this project is the telecommunication

This should be a section of about one page. Here, the context of
the problem should be fully described. Also, previous solutions to the same problem
should be referenced. Most of the bibliography is cited here. While referencing previous solutions, you can also mention computational techniques used, which provides a complete reference for the reader. Usually, this section does not include images or
tables.

\section{Preprocessing data}
Here you should describe the design of your solution, your
technical decisions, and why you believe your solution is appropriate. Some general
images are recommended to enhance your explanation, and this section is expected
to be 1 to 2 pages long. It is not recommended to include code here, but if you are
presenting a complex algorithm, you may add it. Remember to write full paragraphs,
with one idea per paragraph, and avoid using itemized lists. Aim for clarity and
readability. Include enough detail for reproducibility. Briefly mention any limitations
or assumptions in your approach.

\subsection{EDA}
Dataset has the next features:

\begin{table}[H]
\centering
\renewcommand{\arraystretch}{1.3}
\begin{tabular}{|l|p{10cm}|}
\hline
\textbf{Feature} & \textbf{Description} \\ \hline
customerID & Unique value that identify a customer \\ \hline
gender & Determines whether customer is Male or Female \\ \hline
SeniorCitizen & Determines whether a person is citizen or not \\ \hline
Partner & Determines whether person has partner \\ \hline
Dependents & Determines whether customer has a dependent \\ \hline
tenure & Number of months has stayed with the company \\ \hline
PhoneService & Whether customer has phone service \\ \hline
MultipleLines & Whether the customer has multiple lines, or one, or does not have phone service \\ \hline
InternetService & Internet service provider (Fiber optic, DSL or no) \\ \hline
OnlineSecurity & Whether customer has online security, or not, or does not have internet service \\ \hline
OnlineBackup & Whether customer has online backup, or not, or does not have internet service \\ \hline
DeviceProtection & Whether customer has device protection, or not, or does not have internet service \\ \hline
TechSupport & Whether customer has tech support, or not, or does not have internet service \\ \hline
StreamingTV & Whether customer has streaming tv, or not, or does not have internet service \\ \hline
StreamingMovies & Whether customer has streaming movies, or not, or does not have internet service \\ \hline
Contract & Contract term duration (Month-to-month, One year, Two year) \\ \hline
PaperlessBilling & Whether customer has paperless billing or not \\ \hline
PaymentMethod & Payment method (Electronic check, Mailed check, Bank transfer (automatic), Credit card) \\ \hline
MonthlyCharges & Value paid by customer monthly \\ \hline
TotalCharges & Total value paid by customer \\ \hline
Churn & Whether the customer churned or not, stopped using the service last month. \\ \hline
\end{tabular}
\caption{\centering Dataset features description}
\label{tab:features}
\end{table}

\subsection{Subsection 2}
Subsection example 1

\subsection{Subsection 3}
Subsection example 2

\section{Results and Discussion}
Explain the experiments you performed to demonstrate
that your solution works. Summarize your unit tests in a paragraph, mentioning their
philosophy, quantity, and results. Also, discuss any integration and acceptance tests to
address software quality and user expectations. If possible, compare your results with
other solutions using metrics. Use tables to summarize test definitions and results, and
charts (such as box-plots or time series) for comparisons, depending on your project.
Highlight the most important findings. Every figure and table must have a descriptive
caption and be referenced in the text.

\section{Conclusion}
Write a couple of paragraphs to summarize your work, results, and
achievements. Clearly state why your work was successful and how it helped to solve
the problem. End with suggestions for future work or open questions that remain

\begin{thebibliography}{9}
\bibitem{rummery1994}
Rummery, G. A., \& Niranjan, M. (1994). On-line Q-learning using connectionist systems. \textit{Technical Report}, Cambridge University Engineering Department.


\end{thebibliography}

\end{document}